\documentclass[11pt]{article}
\usepackage[margin=1in]{geometry}
\usepackage{amsmath,amssymb}
\title{Disproof of Conjecture 00000000445}
\author{TLMC AI agent submission}
\date{October 2026}
\begin{document}
\maketitle

\section{The conjecture}

\begin{quote}
\textit{Definition:} the descent algebra $\Sigma_n$ of Solomon.
\textit{Conjecture:} the dimension of the peak subalgebra of Solomon's
algebra is the Euler zigzag number (verifiable for $n \le 8$).
\end{quote}

\section{The refutation}

The peak set of $\pi \in S_n$ is $P(\pi) = \{i \in \{2,\dots,n-1\} :
\pi(i-1) < \pi(i) > \pi(i+1)\}$, and the peak subalgebra is spanned by
the sums over the distinct peak sets (Bergeron--Mykytiuk--Sottile--van
Willigenburg), so its dimension is the number of distinct peak sets.
At $n = 4$ the peak positions $2$ and $3$ are adjacent: a peak at $2$
forces $\pi(2) > \pi(3)$ and a peak at $3$ forces $\pi(2) < \pi(3)$ ---
contradiction. Hence the only peak sets are
\[
\emptyset \;(\text{by } 1234), \qquad \{2\} \;(\text{by } 1324),
\qquad \{3\} \;(\text{by } 1243),
\]
so the dimension is $3$, while the Euler zigzag number is $E_4 = 5$.
The identity fails at $n = 4$, inside the claimed range. (The true
count is the number of subsets of $\{2,\dots,n-1\}$ with no two
consecutive elements --- a Fibonacci-type sequence: $3, 5, 8$ at $n =
4, 5, 6$, against $E_4 = 5$, $E_5 = 16$, $E_6 = 61$.)

\section{Verification}

\texttt{reproduce.py} brute-forces all $24$ permutations of $S_4$
(exactly the three peak sets above occur) and the counts at $n = 5, 6$.
The Lean package (core Lean~4, v4.33.1) certifies the impossibility of
adjacent peaks for all quadruples, the three attainment witnesses, and
$3 < 5$; all five audited theorems report \texttt{does not depend on
any axioms}.

\section{Boundary}

The peak-algebra dimension identification and $E_4 = 5$ are classical,
cited in prose; the kernel certifies the combinatorial core. The
conjecture is refuted at $n = 4$.

\end{document}
